%%%%%%%%%%%%%%%%%%%%%%%%%%%%%%%%%%%%%%%%
\chapter{Conclusion and future work}
\label{chap:conclusion}
%%%%%%%%%%%%%%%%%%%%%%%%%%%%%%%%%%%%%%%%

Cette dernière partie ne doit pas être longue : trois pages
suffisent. Il s'agit de rédiger une sorte de bilan de votre
recherche. C'est donc en fin de travail que cette partie sera
écrite. La conclusion a pour objectifs de récapituler les axes
essentiels de votre travail et de mettre en évidence l'apport original
de votre travail. Si des pistes de réflexion n'ont pas pu être suivies
(pour diverses raisons à signaler) mais qu'elles ont un intérêt
particulier, elles peuvent être explicitées dans cette partie. De
plus, les pistes à explorer à la suite à votre travail et les
perspectives qui découlent doivent également figurer dans votre
conclusion. Si vous identifiez des limites et/ou si vous envisageriez
d'autres façons de procéder si c'était à refaire, vous pouvez aussi en
discuter dans votre conclusion.